\section{方法}

\subsection{软件范围与冻结配置}

本文方法使用 Pymatgen 解析晶体结构，并以 NumPy 与 scikit-learn 实现数值和模型流程\cite{Ong2013,Harris2020,Pedregosa2011}。方法对应仓库分支 \code{codex/naconductor-repair} 中的 \code{automat-naconductor} 实现。运行配置由 \code{run\_info.yaml} 提供，关键默认值包括：目标列 \code{log\_sigma}，结构列 \code{cif\_path}，体系列 \code{system}，阴离子列 \code{anion\_type}；Ridge 正则化系数 1.0；稳定性 Lasso 系数 0.05；100 次 50\% 无放回子采样；选择频率阈值 0.60；15 个固定标准正态噪声列；组合公式最少 2 个、最多 3 个描述符；重复子采样 10 次、测试比例 0.2、随机种子 42。

冻结身份由原始 CSV 的规范绝对路径、实际载入的 \code{descriptors/\_\_init\_\_.py} 路径和注册表修订号 \code{structural-registry-v1-2026-08-03} 组成。配置中的 \code{data.raw\_file} 与 \code{shared\_input.raw\_file} 必须解析为同一文件，注册表配置必须指向当前 Python 实际导入的模块。Pipeline 和 Agent 的所有持久化工件都保留该身份。

\subsection{CIF 路径解析与失败关闭预检}

相对 CIF 路径始终相对于原始 CSV 所在目录解析，不回退到进程当前目录。Pipeline 在初始化描述符列和写入 CSV/JSON 之前检查所有路径；Agent 除路径存在外，还逐一调用 Pymatgen \code{CifParser} 解析结构。任一缺失路径或不可解析结构都会中止整批运行。单个描述符计算失败可以记录为 NaN 供审计，但若一个描述符不足 5 个有限值，或 Stage 0 没有任何达到覆盖阈值的真实活动描述符，则分析失败。

\subsection{化学物种、占位和周期几何}

元素身份通过归一化元素符号识别，而不是比较带电物种的显示字符串。Na 总量、阴离子、框架元素和部分占位均按元素占位求和。周期性间隙候选来自实际 Voronoi 顶点；距离使用晶格最小像，并以周期最小像进行候选去重。Wyckoff 多样性定义为 \code{SpacegroupAnalyzer.get\_symmetrized\_structure()} 返回的 \code{equivalent\_indices} 组数。

\subsection{描述符覆盖与噪声列}

批量特征化默认只计算 38 个活动描述符。每列至少需要在 50\% 样本中为有限值，才能进入特征矩阵。描述符原始值与 NaN 原样保留。以固定种子 42 生成 15 列 $\mathcal N(0,1)$ 噪声，列名为 \code{noise\_000} 至 \code{noise\_014}。噪声列用于稳定性基线，不进入物理代表或组合候选。

\subsection{秩感知去混杂}

system 和 anion 均采用带参考类别的虚拟编码。构造带截距的 system 设计矩阵后，按固定列顺序逐一尝试加入 anion 对比；只有使矩阵秩增加的列才进入实际残差化。审计信息包括 system 秩、全部控制的秩、阴离子增量秩、冗余/增量列和最终控制列。

对每个至少有 5 个有限描述符--目标配对的活动描述符，计算原始 Spearman 和 Ridge 残差上的 Spearman。Ridge 的输入为控制矩阵，不对结构描述符作全局标准化。分类规则和体系代理比边界处理与“结果”部分所述一致。分析结果按去混杂 Spearman 绝对值降序排列。

\subsection{子采样 Lasso 与物理族代表}

Stage 2 的真实特征池严格等于 Stage 1 标签为强物理、弱物理或混合信号的活动描述符；固定噪声全部加入 Lasso。每次子采样新建 \code{SimpleImputer(strategy=median)}、\code{StandardScaler} 和 \code{Lasso(max\_iter=20000)}。选择频率为非零系数出现次数除以 100。噪声基线为 15 个噪声频率的第 95 百分位数；真实特征必须同时超过固定阈值和噪声基线。

代表选择只在稳定且超过噪声基线的描述符中进行。每族按去混杂 Spearman 绝对值取前 $k$ 个，二元模式 $k=1$，三元模式 $k=2$。该容量为硬上限。

\subsection{受约束组合公式}

公式值按左结合顺序计算。加法和乘法使用无序组合避免交换律重复；比例根据有向族规则单独生成。分母必须有限且不等于零。加法要求两个已知物理维度相同；未知外部描述符可以保留以支持兼容注册表。候选至少需要 5 个有限公式--目标配对。

二元和三元候选先按去混杂 Spearman 绝对值排序，最多保留 150 个。三元搜索先在同族两个代表之间应用加法或乘法，再与显式相邻族代表结合。候选结构字段用严格 JSON 写入 CSV；读取时校验 2--3 个组分、$n-1$ 个运算符、声明的 \code{n\_components}、兼容字段及 provenance。

\subsection{交叉验证}

每个训练折独立拟合中位数插补、标准化和 Ridge。阴离子分层使用打乱的 \code{StratifiedKFold}，请求 3 折；最小类别不足 2 时跳过，支持 2 折时显式降折。LOSO 使用 \code{LeaveOneGroupOut}，按 system 留出整组。重复子采样使用 \code{StratifiedShuffleSplit}，精确执行 10 次、测试比例 0.2 和种子 42；若任一 system 不足 2 个样本或训练/测试规模不能覆盖全部类别，则跳过。

单折报告测试 Spearman 和 MAE。每个策略的均值只使用有限折。综合分为所有未跳过且均值有限策略的绝对 Spearman 平均值，同时保存可用策略数、是否三种齐全和 skip 原因。

\subsection{V1--V4 探索性证据}

V1 对每个公式分量分别在 system 内置换，保持体系构成和公式结构不变，生成 100 个匹配噪声公式。V2 首先仅在训练折用秩感知控制拟合目标；测试目标残差由训练控制模型产生。随后在训练折用公式值预测目标残差，并在 held-out 样本上汇总 OOF Spearman。若 system 类别支持分层，最多使用 5 折；否则使用打乱 K 折。至少需要 2 个可用折和有限 OOF Spearman 才标记可用。

V3 在每个 system 内计算公式与目标的原始 Spearman。V4 在每个 system 内有放回抽样，默认 500 次，报告 2.5\% 和 97.5\% 分位数。V4 不重新选择描述符或公式，因此不包含选择不确定性。所有区块均返回 \code{status=exploratory}、\code{available}、\code{reason}，整体固定 \code{causal\_claim=false}。

\subsection{Agent 单描述符审计}

Agent CLI 要求 \code{--descriptor-name}，且名称必须是活动描述符。它对冻结 CSV 的所有 CIF 进行存在性和解析性预检，逐结构计算指定描述符，保存失败数量，并使用与 Pipeline 相同的去混杂和 CV 实现。Agent 结果 TSV 含原始/去混杂 Spearman、代理比、标签、各 CV 指标、综合分、数据覆盖和冻结身份；审计 CSV 保存材料、CIF、体系、阴离子、目标和描述符原始值。

\subsection{自动化测试与静态验证}

修复分支记录的验证命令包括 \code{pytest -q}（70 passed）、\code{python -m compileall -q .} 及五个 CLI 的 \code{--help}。额外冒烟测试验证：缺失 CIF 时不存在任何 Pipeline/Agent 结果目录；旧全缺失特征表被拒绝；合成数据可顺序执行 Stage 1--4，并生成真实三元候选、严格 JSON provenance、V1--V4 和策略 skip 状态。\code{git diff --check} 未发现空白错误。

\section{数据与代码可用性}

代码、配置、测试和本手稿位于 \code{northkd/pull\_file} 仓库的 \code{codex/naconductor-repair} 分支。当前修复提交没有补齐原始表引用的 84 个 CIF，也没有重新生成真实研究结果。补齐结构后应从冻结 raw CSV 重新运行特征化，保存结构文件哈希、软件环境、完整阶段工件和运行提交号。

\section{致谢}

致谢信息待补充。

\section{作者贡献}

作者贡献信息待补充。

\section{利益冲突}

作者声明信息待补充。

\section{补充信息}

建议补充材料包含：41 项公开描述符的完整定义与单位、38 项活动描述符清单、规则注册表、控制矩阵秩审计、全部回归测试目录、合成端到端工件示例，以及补齐 CIF 后的真实 Stage 0--4 输出。当前 PDF 不包含未运行的 DFT、NEB、BVSE 或 AIMD 结果。
